\chapter{Mathematical Models in Machine Learning and Computational Foundations}
\label{chap:computational-models}

\FloatBarrier
\section{The Two Pillars of Computational Mathematics in Machine Learning}

The computational treatment of machine learning and large-scale data analysis rests upon two complementary, deeply intertwined mathematical disciplines~\cite{strang2019linear}:
\begin{itemize}
    \item \textbf{Numerical Linear Algebra:} provides the foundational language to represent geometric transformations, project data onto subspaces, compute matrix factorizations (SVD, QR, Cholesky), and quantify floating-point error propagation via condition numbers~\cite{trefethen1997numerical,demmel1997applied}.
    \item \textbf{Mathematical Optimization:} formalizes the decision-making process through constrained and unconstrained formulations, studies analytical optimality conditions (Fermat's stationary rules and Karush-Kuhn-Tucker conditions), investigates Lagrangian duality, and designs scalable iterative descent algorithms~\cite{boyd2004convex,nocedal2006numerical,bazaraa2006nonlinear}.
\end{itemize}

In real-world Data Science, optimization algorithms cannot function without efficient, stable linear algebra routines, as determining descent directions at each step invariably requires solving a linear algebraic system.

\begin{figure}[H]
\centering
\resizebox{0.95\textwidth}{!}{%
\begin{tikzpicture}[font=\small]

    % Left Pillar Card: anchor=east at (-0.3, 0) ensuring a clean channel
    \node[
        draw=blue!25,
        fill=blue!1.5,
        rounded corners=8pt,
        text width=7.1cm,
        inner sep=10pt,
        align=left,
        anchor=east
    ] (left) at (-0.3, 0) {
        {\large\bfseries\color{blue!80!black} Numerical Linear Algebra}\\[2pt]
        {\scriptsize\bfseries\color{blue!60!black}THE COMPUTATIONAL ENGINE}\\[8pt]
        \footnotesize
        \textbf{Geometry and Representation:}\\
        Orthonormal bases, subspaces, orthogonal projectors.\\[5pt]
        \textbf{Matrix Factorizations:}\\
        Singular values (SVD), spectral, QR, Cholesky.\\[5pt]
        \textbf{Numerical Stability and Errors:}\\
        Condition number $\kappa(A)$, backward error propagation.\\[5pt]
        \textbf{Large-Scale Algorithms:}\\
        Sparse matrix solvers, Krylov/Lanczos methods.
    };

    % Right Pillar Card: anchor=west at (+0.3, 0) ensuring a clean 0.6cm separation
    \node[
        draw=red!25,
        fill=red!1.5,
        rounded corners=8pt,
        text width=7.1cm,
        inner sep=10pt,
        align=left,
        anchor=west
    ] (right) at (0.3, 0) {
        {\large\bfseries\color{red!80!black} Mathematical Optimization}\\[2pt]
        {\scriptsize\color{red!60!black}\bfseries THE DECISIONAL FRAMEWORK}\\[8pt]
        \footnotesize
        \textbf{Formulation of Learning:}\\
        Empirical risk, regularization, decision constraints.\\[5pt]
        \textbf{Optimality Conditions:}\\
        Stationary points, global convexity, KKT conditions.\\[5pt]
        \textbf{Duality Theory:}\\
        Lagrangian dual, duality gap, kernel methods.\\[5pt]
        \textbf{Descent and Convergence:}\\
        Gradient, Newton, Quasi-Newton, scaled directions.
    };

    % Top header: Disciplinary Framing anchored with anchor=south above pillars
    \node[
        draw=gray!30,
        fill=gray!4,
        rounded corners=6pt,
        text width=15.3cm,
        inner sep=8pt,
        align=center,
        anchor=south
    ] (top) at ($(left.north)!0.5!(right.north) + (0, 0.7cm)$) {
        {\large\bfseries\color{gray!90!black} FOUNDATIONS OF COMPUTATIONAL MATHEMATICS FOR DATA ANALYSIS}\\[2pt]
        {\footnotesize\color{gray!70!black}The two complementary pillars of modern machine learning}
    };

    % Bottom integration panel anchored with anchor=north below pillars
    \node[
        draw=teal!30,
        fill=teal!2,
        rounded corners=6pt,
        text width=15.3cm,
        inner sep=9pt,
        align=center,
        anchor=north
    ] (bot) at ($(left.south)!0.5!(right.south) + (0, -0.7cm)$) {
        {\bfseries\color{teal!80!black} Operational Confluence Point: The Machine Learning Algorithmic Loop}\\[4pt]
        \parbox{14.8cm}{\footnotesize\centering
        \textbf{Optimization} formulates the problem (\textit{what} to minimize and under what constraints) and generates at each step a directional system $\nabla^2 f(x) \Delta x = -\nabla f(x)$.\\
        \textbf{Numerical linear algebra} provides the efficient and stable algorithms to solve it across millions of parameters (\textit{how} to compute it efficiently).}
    };

    % Accent stripes on pillars
    \draw[line width=4pt, color=blue!70!black, line cap=round] ($(left.north west)+(2pt,-5pt)$) -- ($(left.south west)+(2pt,5pt)$);
    \draw[line width=4pt, color=red!70!black, line cap=round] ($(right.north west)+(2pt,-5pt)$) -- ($(right.south west)+(2pt,5pt)$);

\end{tikzpicture}%
}
\caption{Synergistic conceptual framework: complementary integration of Numerical Linear Algebra and Mathematical Optimization in Machine Learning.}
\label{fig:pillars-architecture}
\end{figure}

\FloatBarrier
\section{The Philosophy of Mathematical Modeling: Brunelleschi and Galileo}

\subsection{The Genesis of Modeling: From Perspective to Computation}
Extracting actionable insight from massive datasets addresses a fundamental epistemological challenge: raw data is unstructured, cognitively overwhelming, and incapable of making autonomous decisions.

\begin{definition}[Computational Model]
\textbf{Learning (machine learning)} consists of extracting structural regularities from empirical observations and synthesizing them into an abstract mathematical artifact (\textbf{model}), computationally tractable and designed to drive predictions or operational decisions.
\end{definition}

This philosophy has deep historical roots:
\begin{itemize}
    \item \textbf{Filippo Brunelleschi and Perspective:} To erect the free-standing dome of Santa Maria del Fiore in Florence without ground-based wooden falsework, Brunelleschi did not proceed by trial and error on site; he built rigorous mathematical models (geometric projections on an optical cone) capable of predicting radial thrust forces.
    \item \textbf{The Economic Advantage of Mathematics:} Constructing physical scale models (such as wind tunnels in aeronautics) incurs enormous material costs. \textbf{Mathematics costs far less than wood or steel}: equations and algorithms allow simulating billions of operational scenarios at negligible cost.
    \item \textbf{Galileo Galilei and Scientific Abstraction:} Modern physics was born when scientists ceased cataloging every accidental empirical detail and instead abstracted parsimonious mathematical laws expressed in algebraic language.
\end{itemize}

\begin{noteblock}{Definition of a Computational Mathematical Model}
A model is not an identical replica of physical reality. It is an \textbf{abstract mathematical formalization} designed to capture essential causal or correlative patterns while filtering out noise, making the problem tractable on digital computers.
\end{noteblock}

\FloatBarrier
\section{Characterization and Taxonomy of Models}

\subsection{The Triangle of Conflicting Requirements}
Designing any computational model involves balancing three inherently competing objectives:
\begin{enumerate}
    \item \textbf{Accuracy (Accurate):} The fidelity with which the model reflects the actual physical dynamics or behavior of the system.
    \item \textbf{Computational Inexpensiveness (Inexpensive):} Algorithmic simplicity, translating to minimal memory footprints and fast execution times ($\BigO{n}$ or quasi-linear complexity).
    \item \textbf{Generality (General):} The versatility to apply across diverse operational environments and domain distributions.
\end{enumerate}

\begin{property}[Model Incompatibility Principle]
It is theoretically and practically impossible to maximize all three objectives simultaneously: hyper-accurate models require vast parameter counts that destroy computational efficiency; parsimonious, rapid models compromise fine-grained accuracy or generality. Model design is inevitably an \textbf{engineering trade-off}.
\end{property}

\begin{figure}[H]
\centering
\resizebox{0.82\textwidth}{!}{%
\begin{tikzpicture}[>=Stealth, scale=1.0]
    % Triangle vertices
    \coordinate (A) at (0, 3.8);
    \coordinate (B) at (-4.2, 0);
    \coordinate (C) at (4.2, 0);

    % Outline and fill
    \draw[thick, draw=blue!70!black] (A) -- (B) -- (C) -- cycle;
    \fill[blue!4] (A) -- (B) -- (C) -- cycle;

    % Labels
    \node[above=4pt, font=\bfseries\normalsize, text=blue!90!black, align=center] at (A) {Accuracy\\\scriptsize (Analytical fidelity)};
    \node[below left=2pt, font=\bfseries\normalsize, text=green!60!black, align=center] at (B) {Computational Inexpensiveness\\\scriptsize ($\BigO{n}$ complexity)};
    \node[below right=2pt, font=\bfseries\normalsize, text=orange!80!black, align=center] at (C) {Generality\\\scriptsize (Domain transferability)};

    % Center point
    \node[circle, fill=red!80!black, inner sep=2.5pt] (P) at (0, 1.4) {};
    \node[below=5pt, font=\footnotesize\bfseries, text=red!80!black, align=center] at (P) {Engineering Trade-off\\(Operational sweet spot)};

    % Arrows
    \draw[<->, dashed, gray!60, thick] (-1.8, 2.0) -- (-0.4, 2.8);
    \draw[<->, dashed, gray!60, thick] (1.8, 2.0) -- (0.4, 2.8);
    \draw[<->, dashed, gray!60, thick] (-2.2, 0.3) -- (2.2, 0.3);
\end{tikzpicture}%
}
\caption{The triangle of conflicting requirements in mathematical modeling: accuracy, computational inexpensiveness, and generality.}
\label{fig:model-triangle}
\end{figure}

\subsection{Analytical Models vs Data-Driven (Synthetic) Models}
Table~\ref{tab:model-comparison} outlines the core differences between the two primary modeling paradigms.

\begin{table}[H]
\centering
\small
\begin{tabularx}{\textwidth}{@{} l Y Y @{}}
\toprule
\textbf{Property} & \textbf{Analytical Approach (White-Box)} & \textbf{Data-Driven Approach (Black-Box)} \\
\midrule
\textbf{Principio Base} & Models each individual physical component and its differential interactions. & Does not model internal mechanics; reproduces observed input-output behavior directly. \\
\addlinespace
\textbf{Struttura Matematica} & Flexible, non-linear, derived from first physical principles and conservation laws. & Standardized, rigid, parametric (e.g., deep neural networks, SVMs, linear hyperplanes). \\
\addlinespace
\textbf{Parameters} & Few, calibrated parameters with precise, interpretable physical meaning. & Massive parameter counts (millions to billions), estimated automatically via optimization. \\
\addlinespace
\textbf{Development Cost} & High: requires deep domain expertise and invasive internal system access. & Low/scalable: requires only historical observation logs and standardized optimization solvers. \\
\bottomrule
\end{tabularx}
\caption{Comparison between analytical white-box modeling and empirical data-driven black-box modeling.}
\label{tab:model-comparison}
\end{table}

\subsection{Fitting vs Machine Learning: Empirical Risk and Generalization}
Every model is an approximation (\textit{``the map is not the territory''}):
\begin{itemize}
    \item \textbf{Fitting (Empirical Risk Minimization):} Identifies parameter vectors that minimize discrepancy purely on known observations (\textit{training set}):
    \begin{equation}
        \min_{\mathbf{w}} \mathcal{R}_{\text{emp}}(\mathbf{w}) = \frac{1}{m} \sum_{i=1}^m \ell(f(\mathbf{x}^i; \mathbf{w}), y^i).
    \end{equation}
    \item \textbf{Learning (Generalization):} Requires $f(\mathbf{x}; \mathbf{w})$ to maintain low error on unseen data drawn from the same underlying distribution (\textit{test set}). A hyper-parametrized model that drives training error to zero often suffers from \textbf{overfitting}, fitting spurious noise rather than true underlying relationships.
\end{itemize}

\FloatBarrier
\section{Four Core Mathematical Models in Machine Learning}

\subsection{Example 1: Linear Least Squares Estimation}
Given $m$ observation pairs $(\mathbf{x}^i, y^i) \in \R^n \times \R$, we postulate a linear input-output mapping:
\begin{equation}
    f(\mathbf{x}) = \mathbf{w}^T \mathbf{x} + w_0 = \sum_{j=1}^n w_j x_j + w_0.
\end{equation}
Absorbing the bias $w_0$ into the augmented design matrix $X$:
\begin{equation}
    \mathbf{w}_{\text{ext}} = \begin{bmatrix} w_0 \\ \mathbf{w} \end{bmatrix} \in \R^{n+1}, \quad
    X = \begin{bmatrix} 1 & (\mathbf{x}^1)^T \\ \vdots & \vdots \\ 1 & (\mathbf{x}^m)^T \end{bmatrix} \in \R^{m \times (n+1)}, \quad
    \mathbf{y} = \begin{bmatrix} y^1 \\ \vdots \\ y^m \end{bmatrix} \in \R^m.
\end{equation}
The linear least squares problem minimizes the Euclidean residual:
\begin{equation}
    \min_{\mathbf{w}} \mathcal{L}(\mathbf{w}) = \|\mathbf{y} - X\mathbf{w}\|_2^2 = \mathbf{y}^T \mathbf{y} - 2\mathbf{w}^T X^T \mathbf{y} + \mathbf{w}^T X^T X \mathbf{w}.
\end{equation}
Setting the gradient $\nabla \mathcal{L}(\mathbf{w})$ to zero yields the stationary condition, known as the \textbf{Gauss Normal Equations}:
\begin{equation}
    \nabla \mathcal{L}(\mathbf{w}) = -2 X^T \mathbf{y} + 2 X^T X \mathbf{w} = \mathbf{0} \implies X^T X \mathbf{w} = X^T \mathbf{y}.
\end{equation}
When $X$ has full column rank ($\rank(X) = n+1$), $X^T X$ is symmetric positive definite and invertible:
\begin{equation}
    \mathbf{w}^* = (X^T X)^{-1} X^T \mathbf{y}.
\end{equation}

\begin{alertblock}{Golden Numerical Rule: Never Compute $(X^T X)^{-1}$ Explicitly}
In scientific computing, \textbf{one must never explicitly form $X^T X$ nor compute its inverse}. Multiplying $X^TX$ squares the condition number ($\kappa(X^T X) = \kappa(X)^2$), drastically magnifying floating-point roundoff errors. The system is solved stably using QR factorization or truncated SVD (backslash operator in GNU Octave: \texttt{w = X \textbackslash\ y}).
\end{alertblock}

\subsection{Example 2: Low-Rank Approximation and Recommender Systems}
Given a partially observed data matrix $M \in \R^{n \times m}$ (e.g., $n$ users and $m$ movies, with $M_{ij}$ indicating ratings), we seek a low-rank factorization $k \ll \min(n, m)$ such that $M \approx AB$:
\begin{equation}
    A \in \R^{n \times k}, \quad B \in \R^{k \times m} \implies \min_{A, B} \|M - AB\|_F^2 = \sum_{i=1}^n \sum_{j=1}^m (M_{ij} - (AB)_{ij})^2.
\end{equation}

\begin{figure}[H]
\centering
\resizebox{0.92\textwidth}{!}{%
\begin{tikzpicture}[>=Stealth]
    % Matrix M
    \node[draw=blue!80!black, fill=blue!10, thick, minimum width=3.4cm, minimum height=4.2cm] (M) at (0, 0) {};
    \node[font=\bfseries\Large] at (0, 0) {$M$};
    \node[above=3pt, font=\small] at (M.north) {$m$ columns};
    \node[left=3pt, font=\small] at (M.west) {$n$ rows};

    \node[font=\Large] at (2.4, 0) {$\approx$};

    % Matrix A
    \node[draw=red!80!black, fill=red!10, thick, minimum width=1.1cm, minimum height=4.2cm] (A) at (3.8, 0) {};
    \node[font=\bfseries\large] at (3.8, 0) {$A$};
    \node[above=3pt, font=\small] at (A.north) {$k$};
    \node[below=3pt, font=\footnotesize, align=center] at (A.south) {Row profiles\\($n \times k$)};

    \node[font=\large] at (4.8, 0) {$\times$};

    % Matrix B
    \node[draw=green!70!black, fill=green!10, thick, minimum width=3.4cm, minimum height=1.1cm] (B) at (7.0, 0) {};
    \node[font=\bfseries\large] at (7.0, 0) {$B$};
    \node[right=3pt, font=\small] at (B.east) {$k$};
    \node[above=3pt, font=\small] at (B.north) {$m$ columns};
    \node[below=3pt, font=\footnotesize, align=center] at (B.south) {Column profiles\\($k \times m$)};
\end{tikzpicture}%
}
\caption{Low-rank matrix factorization: compressing storage from $nm$ entries to $k(n + m)$.}
\label{fig:low-rank-matrix}
\end{figure}

\subsubsection{Memory Compression and Spectral Properties}
Storing $M$ explicitly requires $n \cdot m$ memory entries; storing $A$ and $B$ requires $k(n + m)$. For $n = 10^6$, $m = 10^5$, and rank $k = 10$:
\begin{align}
    \text{Raw storage} &= 10^6 \times 10^5 = 10^{11} \text{ cells}, \\
    \text{Factorized storage} &= 10 \times (10^6 + 10^5) \approx 1.1 \times 10^7 \text{ cells},
\end{align}
achieving a \textbf{99.99\% storage reduction}.

By the Eckart-Young-Mirsky Theorem, the optimal low-rank matrix under the Frobenius norm is given by the truncated Singular Value Decomposition (SVD).

\begin{alertblock}{Computational Warning on Sparse Matrices}
When $M$ is sparse, \textbf{one must never form $M^T M$ or $M M^T$}. Matrix multiplication destroys sparsity, producing a dense matrix that triggers immediate Out-Of-Memory (OOM) failures. One must rely exclusively on iterative methods using matrix-vector products (Lanczos bidiagonalization).
\end{alertblock}

\subsection{Example 3: Support Vector Machines (SVM) --- Margin, Duality, and Kernels}
For binary classification with labels $y^h \in \{-1, +1\}$, a linear classifier makes predictions via $\hat{y} = \operatorname{sign}(\mathbf{w}_+^T \mathbf{x} - w_0)$.

\begin{figure}[H]
\centering
\resizebox{0.82\textwidth}{!}{%
\begin{tikzpicture}[>=Stealth, scale=1.05]
    % Axes
    \draw[->, thick, gray!60] (-0.5, 0) -- (7.5, 0) node[right, black] {$x_1$};
    \draw[->, thick, gray!60] (0, -0.5) -- (0, 6.6) node[above, black] {$x_2$};
    \draw[help lines, color=gray!15, dashed] (0,0) grid (7.0,6.2);

    % Margin corridor (shaded background)
    \fill[teal!6, opacity=0.7] (0.2, 2.4) -- (3.0, 5.2) -- (5.8, 5.6) -- (1.4, 1.2) -- cycle;

    % Upper margin hyperplane (+1): x2 = x1 + 2.2
    \draw[dashed, thick, draw=blue!75!black] (0.2, 2.4) -- (3.0, 5.2) 
        node[above left=3pt, font=\footnotesize\bfseries, blue!85!black, fill=white, inner sep=1.5pt] {$\mathbf{w}_+^T \mathbf{x} - w_0 = +1$};

    % Optimal separating hyperplane (w^T x - w_0 = 0): x2 = x1 + 1.0
    \draw[very thick, draw=black] (0.8, 1.8) -- (4.8, 5.8) 
        node[above right=3pt, font=\footnotesize\bfseries, black, fill=white, inner sep=1.5pt] {$\mathbf{w}_+^T \mathbf{x} - w_0 = 0$};

    % Lower margin hyperplane (-1): x2 = x1 - 0.2
    \draw[dashed, thick, draw=red!75!black] (1.4, 1.2) -- (5.8, 5.6) 
        node[right=4pt, font=\footnotesize\bfseries, red!85!black, fill=white, inner sep=1.5pt] {$\mathbf{w}_+^T \mathbf{x} - w_0 = -1$};

    % Class +1 interior points (blue circles)
    \foreach \coord in {(0.8, 4.6), (1.4, 5.4), (2.0, 5.8), (0.5, 3.2)} {
        \fill[blue!75!black] \coord circle (3.0pt);
    }

    % Class -1 interior points (red squares)
    \foreach \coord in {(3.6, 1.6), (4.6, 2.0), (5.4, 2.4), (4.2, 0.8), (5.6, 1.4)} {
        \fill[red!75!black] \coord +(-3.0pt,-3.0pt) rectangle +(3.0pt,3.0pt);
    }

    % Support Vectors (strictly on supporting hyperplanes)
    % On H+1: (1.6, 3.8) and (2.6, 4.8)
    \fill[blue!75!black] (1.6, 3.8) circle (3.0pt);
    \draw[line width=1.5pt, green!60!black] (1.6, 3.8) circle (6.5pt);

    \fill[blue!75!black] (2.6, 4.8) circle (3.0pt);
    \draw[line width=1.5pt, green!60!black] (2.6, 4.8) circle (6.5pt);

    % On H-1: (2.8, 2.6) and (3.8, 3.6)
    \fill[red!75!black] (2.8, 2.6) +(-3.0pt,-3.0pt) rectangle +(3.0pt,3.0pt);
    \draw[line width=1.5pt, green!60!black] (2.8, 2.6) circle (6.5pt);

    \fill[red!75!black] (3.8, 3.6) +(-3.0pt,-3.0pt) rectangle +(3.0pt,3.0pt);
    \draw[line width=1.5pt, green!60!black] (3.8, 3.6) circle (6.5pt);

    % Support Vector callouts with clean horizontal pointers
    \node[anchor=east, font=\scriptsize\bfseries, green!50!black, fill=white, inner sep=1.5pt] (sv1) at (0.6, 3.8) {Support Vector};
    \draw[->, thick, green!60!black] (sv1.east) -- (1.35, 3.8);

    \node[anchor=west, font=\scriptsize\bfseries, green!50!black, fill=white, inner sep=1.5pt] (sv2) at (4.8, 3.6) {Support Vector};
    \draw[->, thick, green!60!black] (sv2.west) -- (4.05, 3.6);

    % Normal vector w_+ orthogonal to decision boundary
    \draw[->, ultra thick, purple] (4.4, 5.4) -- (3.7, 6.1);
    \node[above=2pt, font=\small\bfseries, purple, fill=white, inner sep=1pt] at (3.7, 6.1) {$\mathbf{w}_+$};
    \draw[purple!80!black, thin] (4.4, 5.4) +(-0.1, 0.1) -- +(-0.0, 0.2) -- +(0.1, 0.1);

    % Margin width (orthogonal distance between H-1 and H+1)
    \draw[<->, line width=1.3pt, purple!85!black] (2.0, 1.8) -- (0.8, 3.0);
    \node[anchor=north east, font=\footnotesize\bfseries, purple!85!black, fill=white, inner sep=2pt] at (1.2, 2.2) {Margin $\delta = \frac{2}{\|\mathbf{w}_+\|_2}$};

    % Legend
    \node[draw=gray!30, fill=white, rounded corners=2pt, font=\scriptsize, anchor=south east] at (7.3, 0.2) {
        \begin{tabular}{cl}
        \tikz\fill[blue!75!black] (0,0) circle (2.5pt); & Class $y = +1$ \\
        \tikz\fill[red!75!black] (-2.5pt,-2.5pt) rectangle (2.5pt,2.5pt); & Class $y = -1$ \\
        \tikz\draw[line width=1.3pt, green!60!black] (0,0) circle (4pt); & Support Vector
        \end{tabular}
    };
\end{tikzpicture}%
}
\caption{Margin geometry in linear Support Vector Machines: maximizing the separation distance between supporting hyperplanes.}
\label{fig:svm-margin}
\end{figure}

Maximizing the geometric margin $2/\|\mathbf{w}_+\|_2$ is equivalent to minimizing $\|\mathbf{w}_+\|_2^2$:
\begin{equation}
    \min_{\mathbf{w}_+, w_0} \frac{1}{2} \|\mathbf{w}_+\|_2^2 \qquad \text{s.t.} \quad y^h (\mathbf{w}_+^T \mathbf{x}^h - w_0) \ge 1 \quad \forall h = 1, \dots, m.
\end{equation}

For non-linearly separable data, slack variables $\xi_h \ge 0$ yield the \textbf{Soft Margin} problem:
\begin{equation}
\begin{aligned}
    \min_{\mathbf{w}_+, w_0, \boldsymbol{\xi}} \quad & \frac{1}{2} \|\mathbf{w}_+\|_2^2 + C \sum_{h=1}^m \xi_h \\
    \text{s.t.} \quad & y^h (\mathbf{w}_+^T \mathbf{x}^h - w_0) \ge 1 - \xi_h, \quad \xi_h \ge 0 \quad \forall h=1, \dots, m.
\end{aligned}
\end{equation}
Parameter $C > 0$ controls the trade-off between margin width and classification violations: $C \to \infty$ forces hard separation, while moderate values allow controlled margin errors to protect test generalization.

\subsubsection{Non-Linear Lifting, Duality, and the Kernel Trick}
When points are not linearly separable in $\R^n$, a non-linear mapping $\boldsymbol{\phi}: \R^n \to \mathcal{H}$ lifts data into a higher-dimensional feature space:
\begin{equation}
    \mathbf{x} \mapsto \boldsymbol{\phi}(\mathbf{x}).
\end{equation}
The Wolfe dual problem is given by:
\begin{equation}
\begin{aligned}
    \max_{\boldsymbol{\alpha}} \quad & \sum_{h=1}^m \alpha_h - \frac{1}{2} \sum_{h=1}^m \sum_{k=1}^m \alpha_h \alpha_k y^h y^k \langle \boldsymbol{\phi}(\mathbf{x}^h), \boldsymbol{\phi}(\mathbf{x}^k) \rangle_{\mathcal{H}} \\
    \text{s.t.} \quad & \sum_{h=1}^m y^h \alpha_h = 0, \quad 0 \le \alpha_h \le C \quad \forall h=1, \dots, m.
\end{aligned}
\end{equation}
Because data points appear solely through inner products, Mercer's Theorem permits evaluating them via a \textbf{kernel function} $K(\mathbf{x}^h, \mathbf{x}^k) = \langle \boldsymbol{\phi}(\mathbf{x}^h), \boldsymbol{\phi}(\mathbf{x}^k) \rangle_{\mathcal{H}}$. For example, the \textbf{Gaussian RBF Kernel}:
\begin{equation}
    K(\mathbf{x}, \mathbf{z}) = \exp\left(-\gamma \|\mathbf{x} - \mathbf{z}\|_2^2\right)
\end{equation}
corresponds to an infinite-dimensional Hilbert space $\mathcal{H}$, evaluated in $\BigO{n}$ time without ever forming $\boldsymbol{\phi}(\mathbf{x})$ explicitly.

\subsection{Example 4: Clustering and the K-Means Algorithm}
Given an unlabelled dataset $X = \{\mathbf{x}^1, \dots, \mathbf{x}^m\} \subset \R^h$, clustering seeks to partition $X$ into $k$ disjoint subsets $\{X_p\}_{p=1}^k$ maximizing intra-cluster cohesion.

\subsubsection{Biconvex Optimization Formulation}
Introducing $k$ cluster centroids $\mathbf{c}^1, \dots, \mathbf{c}^k \in \R^h$ and binary assignment variables $z_{ip} \in \{0, 1\}$:
\begin{equation}
    \min_{\mathbf{c}, Z} \sum_{i=1}^m \sum_{p=1}^k z_{ip} \|\mathbf{c}^p - \mathbf{x}^i\|_2^2 \qquad \text{s.t.} \quad \sum_{p=1}^k z_{ip} = 1 \; \forall i, \quad z_{ip} \in \{0, 1\} \; \forall i,p.
\end{equation}
The problem is non-convex due to binary integrality $z_{ip} \in \{0, 1\}$ and bilinear terms $z_{ip}\mathbf{c}^p$. However, its biconvex structure enables \textbf{Block Gauss-Seidel alternating optimization}:
\begin{enumerate}
    \item \textbf{Fixing $\mathbf{c}$, optimize $Z$:} Each point is assigned to its nearest Euclidean centroid ($z_{i\bar{p}} = 1 \iff \bar{p} = \arg\min_p \|\mathbf{c}^p - \mathbf{x}^i\|_2^2$).
    \item \textbf{Fixing $Z$, optimize $\mathbf{c}$:} Vanishing the convex quadratic gradient with respect to $\mathbf{c}^p$:
    \begin{equation}
        2 \sum_{i \in I(p)} (\mathbf{c}^p - \mathbf{x}^i) = \mathbf{0} \implies (\mathbf{c}^p)^* = \frac{1}{|I(p)|} \sum_{i \in I(p)} \mathbf{x}^i.
    \end{equation}
    The optimal centroid is the \textbf{vector mean (barycenter)} of all points belonging to cluster $p$.
\end{enumerate}

\begin{algorithm}[H]
\caption{K-Means Algorithm for Euclidean Clustering}\label{alg:kmeans-en}
\KwInput{Dataset $X = \{\mathbf{x}^1, \dots, \mathbf{x}^m\} \subset \R^h$, number of clusters $k$, initial centroids $\mathbf{c}^1, \dots, \mathbf{c}^k$, tolerance $\epsilon > 0$}
\KwOutput{Final centroids $\mathbf{c}^1, \dots, \mathbf{c}^k$ and partition $\{I(p)\}_{p=1}^k$}
$v \leftarrow \infty$\;
\Repeat{$(v - v_{\mathrm{new}}) \le \epsilon$}{
    \tcp{Step 1: Point assignment to nearest centroid}
    \For{$p \leftarrow 1$ \KwTo $k$}{
        $I(p) \leftarrow \emptyset$\;
    }
    \For{$i \leftarrow 1$ \KwTo $m$}{
        $p^* \leftarrow \arg\min_{p \in \{1, \dots, k\}} \|\mathbf{c}^p - \mathbf{x}^i\|_2^2$\;
        $I(p^*) \leftarrow I(p^*) \cup \{i\}$\;
    }
    \tcp{Step 2: Centroid re-computation (mean)}
    \For{$p \leftarrow 1$ \KwTo $k$}{
        \If{$I(p) \neq \emptyset$}{
            $\mathbf{c}^p \leftarrow \frac{1}{|I(p)|} \sum_{i \in I(p)} \mathbf{x}^i$\;
        }
    }
    \tcp{Step 3: Objective value evaluation}
    $v_{\mathrm{new}} \leftarrow \sum_{p=1}^k \sum_{i \in I(p)} \|\mathbf{c}^p - \mathbf{x}^i\|_2^2$\;
    $v \leftarrow v_{\mathrm{new}}$\;
}
\Return{$\mathbf{c}^1, \dots, \mathbf{c}^k$}\;
\end{algorithm}

\begin{noteblock}{Why Scientific Computing Violates Agile Modularity}
In standard software engineering (Agile), systems are decomposed into decoupled, independently swappable modules. In scientific computing, this separation fails:
\begin{itemize}
    \item Replacing the Euclidean metric $\ell_2$ with the Manhattan distance $\ell_1$ completely changes the closed-form centroid: the $\ell_1$ optimum is no longer the arithmetic mean, but the \textbf{component-wise median}.
    \item The tight interdependence between loss formulation and update routines demands an integrated, holistic mathematical design.
\end{itemize}
\end{noteblock}

\FloatBarrier
\section{Methodological Principles of Computational Learning}

\subsection{The ``Magic Academy'' Metaphor}
To illustrate the difference between superficial library usage and rigorous computational competence, consider the \textbf{Magic Academy}:
\begin{itemize}
    \item \textbf{The Sorcerer's Apprentice (\textit{Black-Box Practitioner}):} Invokes library routines (\texttt{import sklearn}, \texttt{fit()}) as pre-packaged incantations. In nominal conditions, output appears satisfactory; however, when confronted with numerical instability or ill-conditioning, they lack the insight required to diagnose the root cause.
    \item \textbf{The Master Magician (\textit{Computational Scientist}):} Understands spectral properties and convergence theory, detects ill-posed systems, tracks floating-point error amplification, and replaces fragile algorithms with robust, numerically sound alternatives.
\end{itemize}

\subsection{The Three Guiding Principles of Machine Learning}
\begin{enumerate}
    \item \textbf{Learning is Numerical Computation:} Every AI algorithm fundamentally reduces to solving a linear algebraic system or minimizing an objective loss function.
    \item \textbf{The Dimensional Scale Challenge:} Machine learning problems are rarely mathematically intractable; they are predominantly linear, convex, or quadratic. \textbf{Their difficulty stems entirely from their colossal scale} (\textit{«hard because large, not hard because hard»}), spanning millions to billions of parameters.
    \item \textbf{Linear Algebra as the Engine of Optimization:} Optimizing at extreme scales demands ruthless hardware efficiency: numerical linear algebra provides the computational engine that makes optimization fast, scalable, and numerically stable.
\end{enumerate}

\FloatBarrier
\section{Synoptic Formula Reference of Introductory Models}

Table~\ref{tab:model-summary-en} summarizes the foundational mathematical formulations and their computational paradigms.

\begin{table}[H]
\centering
\footnotesize
\begin{tabularx}{\textwidth}{@{} >{\bfseries}l >{\raggedright\arraybackslash}p{3.8cm} >{\raggedright\arraybackslash}p{3.1cm} Y @{}}
\toprule
\normalfont\textbf{Model} & \textbf{Formulation} & \textbf{Paradigm} & \textbf{Computational Notes} \\
\midrule
Least Squares & $\min_{\mathbf{w}} \|\mathbf{y} - X\mathbf{w}\|_2^2$ & Normal equations:\newline $\mathbf{w}^* = (X^TX)^{-1}X^T\mathbf{y}$ & Solve via QR factorization; never invert $X^TX$ explicitly. \\
\addlinespace[2pt]
Low-Rank & $\min_{A, B} \|M - AB\|_F^2$ & Truncated SVD (Eckart-Young) & Compression $nm \to k(n+m)$; avoid Gramian $M^TM$ on sparse data. \\
\addlinespace[2pt]
SVM (Hard) & $\min \frac{1}{2}\|\mathbf{w}_+\|_2^2$\newline $\text{s.t. } y^h(\mathbf{w}_+^T\mathbf{x}^h - w_0) \ge 1$ & Convex Quadratic Program (QP) & Maximizes geometric margin $2/\|\mathbf{w}_+\|_2$ under strict separation. \\
\addlinespace[2pt]
SVM (Soft) & $\min \frac{1}{2}\|\mathbf{w}_+\|_2^2 + C \sum \xi_h$\newline $\text{s.t. } y^h(\mathbf{w}_+^T\mathbf{x}^h - w_0) \ge 1-\xi_h$ & QP with linear constraints (equivalent to Hinge loss) & Trade-off parameter $C > 0$ balances bias-variance via slack variables. \\
\addlinespace[2pt]
SVM (Kernel) & $\max_{\boldsymbol{\alpha}} \mathbf{1}^T\boldsymbol{\alpha} - \frac{1}{2}\boldsymbol{\alpha}^T \tilde{K} \boldsymbol{\alpha}$\newline $\text{s.t. } \mathbf{y}^T\boldsymbol{\alpha}=0, \; 0 \le \alpha_h \le C$ & QP with box constraints & Inner products in infinite-dimensional Hilbert spaces (RBF) in $\BigO{n}$. \\
\addlinespace[2pt]
K-Means & $\min \sum_{i,p} z_{ip} \|\mathbf{c}^p - \mathbf{x}^i\|_2^2$\newline $\text{s.t. } \sum_p z_{ip}=1, \; z_{ip}\in\{0,1\}$ & Alternating minimization (Block Gauss-Seidel) & Finite convergence to local minimum; cluster centroids given in closed form. \\
\bottomrule
\end{tabularx}
\caption{Synoptic reference of introductory mathematical models and computational properties.}
\label{tab:model-summary-en}
\end{table}
